

\subsection*{Equação Paramétrica da Reta}


\begin{frame}[label=retas]{Equação Paramétrica de uma Reta}

Fixados um ponto ${\color{red} A}$ de uma reta $r$ e um vetor ${\color{blue}\vec{v}}$ paralelo a esta reta, podemos descrever os pontos $P$ dessa reta da seguinte forma:
\[P={\color{red}A}+t{\color{blue}\vec{v}},\ t\in \R.\]
	\begin{center}
		\begin{tikzpicture}
		\tikzset{>=latex}
		\draw[thick] (-1,-1) -- (2,2);
		\node[left] at (2,2) {r};
		\draw[fill=red] (0,0) circle [radius=1.5pt];
		\node[above left] at (0,0) {$\textcolor{red}{A}$};
		\draw[thick, blue,->] (2,0) -- (0.7+2,0.7);
		\node[blue, right] at (0.7+2,0.7) {$\vec{v}$};
		\onslide<2->{
		\draw[fill=black] (1.5,1.5) circle [radius=1.5pt];
		\node[above left] at (1.5,1.5) {$P$};}
	
	\onslide<3->{
\draw[thick, blue,->] (0,0) -- (1.5,1.5);	
  }
		\end{tikzpicture}
	\end{center}
	
Esta equação é dita {\color{blue} equação vetorial paramétrica da reta $r$}, $t$ é dito ser um {\color{blue}parâmetro} e ${\color{blue}\vec{v}}$ é chamado de {\color{blue}vetor diretor.}

\end{frame}

\begin{frame}[label=retas]
Se	$P=(x,y)$, $\textcolor{red}{A=(x_0,y_0)}$ e $\textcolor{blue}{\vt{v}=(a,b)}$
em coordenadas temos que
\[	\left\{\begin{array}{l}
	x=\textcolor{red}{x_0}+\textcolor{blue}{a}t\\
	y=\textcolor{red}{y_0}+\textcolor{blue}{b}t,\  t\in \R.
	\end{array}\right.\]
ou
\[
\begin{bmatrix} x \\ y\end{bmatrix}
={\color{red}\begin{bmatrix} x_0 \\ y_0 \end{bmatrix}}
+t{\color{blue}\begin{bmatrix} a \\ b \end{bmatrix}}
\]
ou simplesmente
\[
(x,y)={\color{red}(x_0,y_0)}+t{\color{blue}(a,b)}.
\]

\begin{exe}\label{exer4} \noindent
 Determine as equações paramétricas da reta que passa pelos pontos  $A=(-1,0)$ e $B=(0,1)$.
\end{exe}
\end{frame}
